% ============================================================================
% CH 9 — HIDDEN EFFECTS                                [slug: hidden-effects]
% STATUS: mechanism inventory rich; measurements pending (login-node
% instruments + SC runs). The pedagogical "challenge ladder" chapter
% (defined 2026-07-12_1224): algorithm first, then confront the reader with
% concealed mechanisms one by one. Granularity flexible — could compress to a
% section inside the merge chapter if thin; currently material justifies a chapter.
% STRUCTURE DELIBERATELY UNFORMED (Matteo, 2026-07-14, outline-refinement
% walk): this chapter was consciously LEFT UNTOUCHED — it is too far in the
% future to shape. Expect better groupings of effects once measured, and
% expect some rungs to be DISCARDED as hypothesized-but-never-measured/
% confirmed. The section set will change a ton; do not restructure or polish
% headers here until the experiments have run. Current headers = inventory
% placeholders only.
% Material map: ../outline.md#hidden-effects
% ============================================================================
\cleardoublepage
\chapter{Hidden Effects: What the Sums Conceal (6--8)}   % page goal — scaffolding, DELETE before submission
\label{ch:hidden}
\vault{2026-07-12_1224_length-channel-refund-hypothesis-hidden-effects-section}

\section{Reading a Composed Model Honestly}
\label{sec:hidden-intro}
% Chapter design: each rung names a mechanism invisible in the model equation,
% derives it, then measures it. Scrutiny-phase findings get dedicated
% measurement runs before any number is quoted here.

\section{The Implicit Item Intercept}
\label{sec:hidden-bt}
\vault{2026-07-12_1808_bt-implicit-item-intercept-shift-asymmetry-directives, 2026-07-12_1815_ui-merge-intercept-feature-attributable-check-which-side-is-free}
% B(t) = 1^T t from the shifted cosine; U/I asymmetry (rank-active vs
% rank-inert). DIRECTIVES: describe, measure (variance share + popularity
% correlation), display (first-class renderer line).

\section{Share Identifiability and Feature Twins}
\label{sec:hidden-shares}
\vault{2026-07-11_2310_dc01-share-identifiability-problem-and-thesis-framing}
% F-DC01-06: correlated features make per-feature splits init-arbitrary;
% grouped-share quotient + determinacy annotations; twins as the extreme case.

\section{The ID--Twins--Popularity Braid}
\label{sec:hidden-braid}
\vault{2026-07-12_0006_id-twins-popularity-braid-cold-eval, 2026-07-11_1556_bias-channel-popularity-prototype-interpretability, 2026-06-16_2014_popularity-bias-prototype-interpretability-ablation}
% How ID rows, twin items, and popularity entangle under cold evaluation;
% registered hypothesis: warm feature-explained fraction predicts cold
% degradation. (User: "insights we absolutely should report.")

\section{Token Mass and Popularity Coupling}
\label{sec:hidden-token-mass}
\vault{2026-07-12_2208_ml1m-token-mass-raw-bags-popularity-coupling, 2026-07-12_2313_ml1m-popularity-coupling-live-magnitude}
% ml-1m raw indicator bags: popular movies speak ~16:1 louder; headwind for
% individuation; live magnitudes.

\section{The Length Channel and Its Refund}
\label{sec:hidden-length-refund}
\vault{2026-07-12_1224_length-channel-refund-hypothesis-hidden-effects-section}
% Registered hypothesis: angular twin tax paid at fI/fU, refunded at the fUfI
% merge on twins/popular items. [Test lands with Ch 8 runs.]

\section{The ids-Arm Dominance Theorem}
\label{sec:hidden-ids-dominance}
\vault{2026-07-12_1739_ids-arm-dominance-theorem-thesis-placement, 2026-07-12_1830_m5-noise-for-crowding-ids-division-of-labor-centering-option}
% DIRECTIVE: goes into the thesis, meticulously checked (itemized obligations):
% feature structure with ID kept costs nothing on the training objective;
% generalization boundary. Anchors the ids-vs-noid discussion.

\section{How Reliable Is the Cosine Read-Out?}
\label{sec:hidden-cosine}
\vault{2026-06-16_2102_cosine-readout-caveat-steck}
% Steck caveat inherited hard (ProtoMF-is-MF); computation fidelity intact,
% semantic read-out at risk; reliability probe.
